\section{Introduction}
\label{sec:introduction}

\subsection{Motivation}
\label{sec:intro-motivation}

A newly launched exchange, referred as exchange
$A$, has mandated our firm as the sole Designated Market Maker for its
EUR/USD spot book. Two established venues, $B$ and $C$, remain the primary
sites of price discovery, carrying on average $75\%$ and $25\%$ of the daily
volume respectively. We receive market data from $B$ with a $200\,$ms delay
and from $C$ with a $170\,$ms delay, while high-frequency participants reach
the same data in $50\,$ms and can act on exchange $A$ before our own view of
fair value has refreshed. The mandate commits us to continuously quoting ten
levels on each side of the book at a common $1\,$bp tick, with a starting
capital $K$ split evenly between EUR and USD and an asymmetric fee grid. A delta limit caps the EUR inventory we are willing
to warehouse before externalising risk as taker flow on $B$ and $C$.

Three constraints shape every design decision that follows. First, on the
day exchange $A$ goes live no historical fills or counterparty behaviour are
observable, so the initial quoting policy cannot be learnt from data and
must be derived from first-principles microstructure models. Second, the
latency gap to HFTs exposes every resting quote to adverse selection, which
the pricing model must compensate for rather than ignore. Third, because we
cannot market-make on $B$ or $C$, the delta limit turns inventory management
into an active cross-venue hedging problem rather than a passive
book-balancing exercise.

\subsection{Objective}
\label{sec:intro-objective}

Our goal is to deliver an end-to-end market-making system that addresses the
three constraints above and can be evaluated under the conditions described
in the subject brief. The project is split into three phases matching the
mandate timeline. Phase 1 bootstraps liquidity without any trading history,
using a heuristic quoting policy derived from classical market-making
theory. The fair mid is reconstructed from the two reference venues at
their respective latencies, the bid-ask spread is sized to compensate for
both inventory risk and adverse selection, and quotes are skewed in the
direction of recent order-flow imbalance. A dedicated hedge router
externalises excess EUR inventory on $B$ and $C$ whenever the delta limit
is approached. Phase 2 re-estimates the Phase 1 parameters by maximum likelihood
from a simulated month of fills. Phase 3 studies how our quoting policy
should degrade from a price-competitive quoter to a reliable fallback
liquidity provider once aggressive HFT market makers populate the inside of
the book.

\subsection{Simulation Environment}
\label{sec:intro-simulation}

Because the exchange has no history at launch and we cannot run the strategy
against real counterparties, the evaluation environment is a layered
simulation built specifically for this project. A stochastic mid-price
generator (GBM, GARCH with jumps, or Heston) produces the fair value
underlying both reference venues. A venue simulator wraps that path into a
realistic bid/ask and depth profile per venue, with a choice of five spread
regimes (static, stochastic, adaptive, asymmetric, and skewed)
reflecting the microstructure assumptions we want to stress. A
limit order book with a matching engine and fill events sits at the centre
of exchange $A$; a Poisson client-flow generator with
truncated Pareto sizes delivers counterparty orders whose intensity decays
exponentially in the quoted half-spread; and a reporting layer reconstructs
the full mark-to-market and realised P\&L decomposition from the trade
history.

\subsection{Evaluation}
\label{sec:intro-evaluation}

Each phase is evaluated on the subject-mandated one-month simulated
backtest, with the capital $K$ reset to its initial split at the start of
every phase. For Phase 1 we report realised and mark-to-market P\&L, the
percentile envelope of mark-to-market as a function of time since trade
inception, top-ten trades by size, full-versus-partial fill rates at each of
the ten quoting levels, and inventory trajectories against the delta limit.
For Phase 2 we compare the heuristic and calibrated configurations on an
independent month using the same diagnostics, so that any difference in
performance is attributable to the updated parameters rather than to the
market regime. For Phase 3 we measure both our own P\&L under HFT
competition and the quality of liquidity observed by exchange-$A$ clients,
distinguishing small from large trades and quantifying the continuity of
two-sided quoting through volatility spikes, one-sided HFT retreats, and
venue outages.

\subsection{Report Structure}
\label{sec:intro-structure}

The remainder of the report follows the step-by-step structure of the
subject. \autoref{sec:orderbook} introduces the order-book data structure
and matching rule; \autoref{sec:market-simulator} describes the venue
simulator and the family of spread regimes supported;
\autoref{sec:quoter} presents the quoting framework in full;
\autoref{sec:pnl} details the inventory and P\&L bookkeeping; \autoref{sec:backtest} gathers the Phase 1 backtest outputs;
\autoref{sec:phase2} analyses the calibration results;
\autoref{sec:phase3} examines the Phase 3 fallback behaviour; and
\autoref{sec:conclusion} summarises the findings and outlines the
remaining modelling questions.
